
\section{The Origins of Zero-Knowledge}

Zero-Knowledge proofs were invented in \textbf{1985} by Shafi Goldwasser, Silvio Micali, and Charles Rackoff at MIT. Their paper --- \textit{The Knowledge Complexity of Interactive Proof Systems} --- introduced the concept of an interactive proof where a prover convinces a verifier of the truth of a statement while revealing nothing beyond that truth. The paper won the Gödel Prize in 1993 and is widely considered one of the most influential works in the history of cryptography.

At the time, the result was purely theoretical. The proof systems described required many rounds of interaction, were expensive to compute, and existed only to answer a complexity-theoretic question: \emph{how much ``knowledge'' does a proof leak?} The answer, in some cases, was: none.

\section{From Theory to Practice}

\subsection*{The Schnorr Protocol}

A landmark step toward practicality came with Claus-Peter Schnorr's identification scheme (1989). The \textbf{Schnorr protocol} is a three-move sigma protocol:

\begin{enumerate}
  \item The prover commits to a random value.
  \item The verifier sends a random challenge.
  \item The prover responds with a value that combines the challenge with their secret.
\end{enumerate}

The verifier can check the response without learning the secret. Sigma protocols like this are the ancestors of almost every practical ZK system in use today.

\subsection*{The Fiat--Shamir Heuristic}

In 1986, Amos Fiat and Adi Shamir showed how to convert any sigma protocol into a \textbf{non-interactive} proof by replacing the verifier's random challenge with a hash of the prover's commitment. The resulting proof can be generated offline and verified by anyone --- no live interaction needed. This single insight is what makes ZK proofs viable on blockchains.

\subsection*{Schnorr Signatures on Cardano with FROST}

Schnorr signatures eventually became a production cryptographic primitive. Today you can use them on Cardano through the \textbf{FROST} threshold signature protocol, accessible at \texttt{https://www.frost.express/}. FROST allows a group of participants to collaboratively sign a message with a Schnorr signature in a way that is indistinguishable from a single-signer scheme on-chain, with no single party holding the full private key.

\section{The SNARK Revolution}

Through the 2000s and early 2010s, researchers worked to make ZK proofs efficient enough for real applications. The breakthrough came in 2013--2016 with a family of constructions called \textbf{zk-SNARKs}.

\paragraph{Trusted setup and Groth16.}
The most efficient SNARK construction is \textbf{Groth16} (Jens Groth, 2016). It produces proofs of only $\sim$200 bytes that verify in milliseconds --- fast enough for on-chain use. The catch is a \emph{trusted setup ceremony}: a group of participants must generate public parameters, and if any participant keeps their secret ``toxic waste'', they could forge proofs. In practice, multi-party ceremonies with hundreds of participants make this risk negligible.

\paragraph{Zcash and real-world deployment.}
\textbf{Zcash} (2016) was the first large-scale deployment of zk-SNARKs. It uses Groth16 to let users make shielded transactions: the blockchain verifies a proof that the sender has enough funds and the transaction is balanced, without revealing amounts or addresses.

\paragraph{Universal setups.}
Circuit-specific trusted setups are impractical for general use. \textbf{PLONK} (2019) and \textbf{Marlin} introduced \emph{universal} setups: a single ceremony valid for any circuit up to a given size. This enabled the modern ecosystem of general-purpose zkVMs and zkEVMs.

\section{Modern ZK: STARKs and Beyond}

\paragraph{STARKs.}
\textbf{Eli Ben-Sasson} and colleagues at the Technion introduced \textbf{zk-STARKs} in 2018. STARKs require \emph{no trusted setup} --- all parameters are public and derived from collision-resistant hash functions. They are also believed to be post-quantum secure. The trade-off is proof size: a STARK proof can be tens to hundreds of kilobytes versus a few hundred bytes for a SNARK.

\textbf{StarkWare} built an entire ecosystem around STARKs: \textit{StarkEx} powers major DeFi protocols, and \textit{StarkNet} is a permissionless L2 on Ethereum.

\paragraph{ZK Rollups.}
The rollup model took off around 2020--2022. \textbf{zkSync}, \textbf{Scroll}, and \textbf{Polygon zkEVM} all use ZK proofs to compress thousands of Ethereum transactions into a single on-chain verification, achieving Ethereum security with a fraction of the gas cost.

\paragraph{Recent developments.}
The frontier is moving fast: \emph{folding schemes} like \textbf{Nova} (2021) allow recursive proof composition with near-zero overhead; \emph{lookup arguments} like Plookup dramatically reduce the cost of proving table lookups and bit operations; and \emph{zkVMs} such as RISC Zero let developers write programs in Rust and get a ZK proof of correct execution essentially for free.

\section{ZK Arrives on Cardano}

\subsection*{Plutus V3 and BLS12-381 Primitives}

With the Chang hard fork (2024), Cardano gained \textbf{Plutus V3} and a set of native BLS12-381 elliptic curve primitives. These allow on-chain verification of Groth16 proofs without any workarounds. The primitives are available in Aiken's standard library:

\begin{itemize}
  \item \texttt{aiken/crypto/bls12\_381/g1} --- operations on the $\mathbb{G}_1$ group
  \item \texttt{aiken/crypto/bls12\_381/g2} --- operations on the $\mathbb{G}_2$ group
  \item Pairing operations: \texttt{miller\_loop} and \texttt{final\_verify}
\end{itemize}

See the full API at \texttt{https://aiken-lang.github.io/stdlib/aiken/crypto/bls12\_381/g1.html}.

\subsection*{Janus Wallet}

The most complete open-source ZK project on Cardano today is \textbf{Janus Wallet}, available at:

\begin{center}
  \texttt{https://github.com/leobel/janus-wallet}
\end{center}

Janus demonstrates a working Groth16 verifier on Cardano: a user generates a ZK proof off-chain that they know a secret, and the on-chain Aiken validator verifies the proof and releases the UTxO. It is the closest thing Cardano has to a production-grade ZK smart contract today.

\subsection*{What This Book Covers}

The rest of this book is structured as a progressive journey:

\begin{enumerate}
  \item \textbf{Starting from basics} --- proving knowledge of a secret value and using a ZK proof to spend a UTxO on Cardano.
  \item \textbf{Compression} --- using off-chain computation with on-chain verification to batch many operations into a single transaction.
  \item \textbf{Advanced constructions} --- privacy-preserving patterns studied for academic purposes. All material here is for educational use. This book does not promote protocols designed to obscure transaction data.
\end{enumerate}
